% Standalone insertion draft: review and integrate into ESM_1.tex after results are final.
% No uncompleted GPU result is asserted here.
\subsection{Automatic coefficient calibration and transfer analysis}

To investigate dependence on dataset-specific manual coefficients, we distinguish
calibration of the original support formula from a complementary calibration
baseline with shared feature definitions. The original formula retains its
dataset-specific evidence features, while a fixed validation search selects the
shortcut penalty, contradiction penalty, and rank-prior mixture from a declared
grid of 224 configurations per dataset. Selection uses validation MRR@20 only;
the selected configuration is frozen before test evaluation.

The shared-feature baseline uses the original rank prior, source-subgraph contact,
bounded log degree, bounded log shared-neighbor count, and a negative direct-contact
indicator. These features are computed without reference to evaluation labels and
use the same definitions and normalization constants across datasets. They are
topological descriptors rather than annotations of therapeutic mechanism. An
equal-weight baseline is included to separate effects of the shared features from
effects of coefficient learning.

For learned calibration, a linear candidate score is fitted using a pairwise
logistic objective with L2 regularization. The primary formulation constrains
coefficients to be nonnegative because feature signs are specified explicitly;
an unconstrained counterpart is also evaluated. The regularization coefficient
is selected by five-fold query-group cross-validation within the existing
validation split, which serves as a calibration pool. Candidates known to be
other positives for the same calibration-training query are excluded as negative
pairs. Queries whose gold entity is absent from the fixed top-20 list cannot
contribute a positive training pair, but remain in all evaluation denominators.

We additionally perform leave-one-dataset-out coefficient transfer: the target
dataset contributes no validation labels to fitting or hyperparameter selection.
Coefficients are fitted on the other datasets, with equal total pair weight per
source dataset, and regularization is selected using leave-one-source-dataset-out
validation. The target graph, structural retriever, and prepared subgraphs remain
available. This experiment therefore assesses transfer of calibration coefficients,
not zero-shot completion on an unseen knowledge graph. Pairwise source-to-target
transfer, feature ablation, calibration-label learning curves, and coefficient
perturbation analyses provide complementary diagnostics.

For fuzzy retrieval, a common edge feature set uses rank-band confidence,
candidate incidence, query incidence, direct-contact penalties, and bounded
degree information, without dataset-specific relation-family priorities.
Automatic selection searches a fixed set of coefficient vectors using validation
generated-answer accuracy, with adjusted MRR@20 as a secondary criterion. Candidate
order and per-query edge budgets are held fixed. A source-only macro-validation
criterion provides the corresponding leave-one-dataset-out edge-coefficient
selection. All selected settings are recorded before test inference.
The source-only restriction applies to this new coefficient-selection step;
the original fuzzy candidate order and per-query budgets are inherited from the
historical dataset-specific pipeline. Thus the complete graph-input pipeline is
not claimed to be free of prior target-dataset calibration.

End-to-end comparisons use one frozen checkpoint per dataset. Candidate-order
interventions retain the full source graph, whereas edge-selection interventions
retain the original fuzzy input candidate order. Exact generated-answer accuracy,
invalid-output rate, and top-candidate-copy rate are reported separately from
adjusted MRR@20. Paired query-cluster bootstrap intervals quantify uncertainty;
multiple gold rows for the same query are resampled together. These analyses are
retrospective checks on previously studied benchmarks, and their exploratory
intervals are not adjusted for multiple comparisons.

\paragraph{Numerical precision correction.}
Following the nonfinite graph-output audit in the information-retention study,
all Comment 2 inference is repeated with graph-encoder checkpoint weights loaded
directly in FP32 and base/PEFT LLM weights loaded directly in BF16. Projected
embeddings must remain finite both before and after the BF16 cast. No clipping,
sample exclusion, or reuse of legacy FP16 predictions is permitted. Candidate
rankings, shared-feature coefficients, and other CPU artifacts remain unchanged.
Every validation edge-vector trial is reevaluated; automatic and source-only
edge selections are frozen again before test inference, using the same search
space, metric definitions, and tie rule.

Identical stable predictions from the Comment 3 run can be reused only when the
full checkpoint, model, precision, decoding, and code fingerprint matches, and
the exact prompt, candidate ordering, and edge multiset match. This is inference
memoization; validation selection does not access test performance. A separate
audit compares FP16 encoding, original-checkpoint FP32 encoding, and casts to
FP16/BF16 on the actual validation/test graph inputs of both runs. Old--new
prediction comparisons identify changed graph inputs after reselection, since
these cases cannot isolate precision effects at a fixed input. The precision
policy changes both graph and LLM arithmetic, so differences are not attributed
causally to either component alone.
For the two adaptive graph arms, an additional diagnostic evaluates the old
selected coefficients at stable precision. This permits an additive metric
decomposition into the precision-policy effect at the fixed old choice and the
effect of validation reselection under stable precision. These ten diagnostic
dataset--variant combinations are evaluated after the main selections are
frozen and do not participate in model selection.
